% STATUS: R1 revised: axis-table layout/citations and comparison sentence split.
% Scope (D9 + Section Blueprint): four threads --- classical planning/search,
% neural algorithm execution, LLM/VLM planning, and runtime verification/tool
% use --- plus a compact axis table (algorithm conditioning, model-owned
% exploration, no-repair runtime, invariant-checked episodes, matched
% modalities, transfer). New citations in iclr2026_conference_related_refs.bib.
\section{Related Work}
\label{sec:related}

% Paragraph purpose: Thread 1 --- classical planning/search and LLM benchmarks.
\paragraph{Classical planning and search as LLM benchmarks.}
A* search~\citep{hart1968formal} and its descendants define the family of heuristic best-first algorithms whose execution semantics our policy learns.
Best-First Width Search (BFWS) combines goal-count exploitation with width-based exploration over a novelty measure~\citep{lipovetzky2017bestfirst} and supplies the structural gate in our evaluation (Section~\ref{sec:design}).
Recent benchmarks have tested whether LLMs can generate plans or reason about change in formal planning domains.
PlanBench~\citep{valmeekam2023planbench} found LLM plan generation far below classical solvers.
A broader critical investigation showed that LLMs fail to plan even in simple domains where chain-of-thought prompting should suffice~\citep{valmeekam2023critical}.
AutoPlanBench~\citep{stein2025autoplanbench} and ACPBench Hard~\citep{kokel2026acpbenchhard} extend these probes with automated prompt generation and harder reasoning tasks, confirming the gap.
The observational study that motivates our intervention~\citep{huang2026twodistinct} diagnosed two non-compensatory competencies: operational reasoning improved with scale and prompting, but structural enumeration did not.
All of these evaluate plans or plan-relevant judgments as end products.
None asks the model to execute a declared search algorithm step by step.

% Paragraph purpose: Thread 2 --- neural algorithm execution.
\paragraph{Neural algorithm execution.}
Searchformer~\citep{lehnert2024beyond} is the nearest prior work in this thread.
It trains a transformer to emit A* search dynamics as serialized state add/remove tokens, then bootstraps to shorter-than-A* traces on Sokoban.
Searchformer operates on a single algorithm and domain family, serializes trace tokens without a typed-operation interface, and has no mechanism to validate or reject emissions during an episode.
Stream of Search~\citep{gandhi2024stream} serializes symbolic search processes, including backtracking and fruitless paths, into a unified language.
A transformer pretrained on solver streams outperforms training on optimal trajectories alone.
The model internalizes the full search process, but training and evaluation provide no per-operation validation or invariant checking.
Earlier work established the broader thread.
\citet{velickovic2020neural} trained GNNs with step-level imitation on individual BFS and Bellman--Ford steps, showing positive transfer across algorithms, though the setup used teacher-shaped per-node targets on synthetic graph encodings with no external runtime.
\citet{zaremba2014learning} trained LSTMs to predict the output of short programs, showing that curriculum design was critical for length generalization.
Our work differs from this thread in its typed operations, whose explicit operands fully determine exploration.
A runtime validates each operation against algorithm-specific invariants during the episode.
The evaluation reports operation validity, invariant compliance, and episode success as separate measurements.

% Paragraph purpose: Thread 3 --- LLM/VLM planning and step/process supervision.
\paragraph{LLM/VLM planning and step-level supervision.}
Several frameworks have the LLM participate in search by contributing local proposals while an external scaffold controls exploration.
Tree of Thoughts~\citep{yao2023tree} scaffolds tree-structured exploration with LM-based evaluation and backtracking, but the expansion and pruning decisions belong to the scaffold, not the model.
RAP~\citep{hao2023reasoning} places the LLM inside Monte Carlo Tree Search as both world model and reasoning agent, yet the UCT policy owns exploration.
In our formulation, the model itself issues the operations that determine which states are expanded and which successors are generated.

Step-level and process supervision have shown advantages over outcome-only training.
Scratchpads~\citep{nye2021show} train transformers to emit intermediate computation steps for algorithmic tasks, improving accuracy on long addition and program execution.
Let's Verify Step by Step~\citep{lightman2023lets} demonstrated that per-step human labels outperform outcome supervision for mathematical reasoning at scale.
These results support our choice of per-decision teacher labels.
The distinction is that scratchpad tokens and human correctness judgments are unchecked free-form annotations.
Our operations are typed and runtime-validated, bound to replay-verified positions in the search trace.

Recent work uses visual intermediate states in VLM reasoning.
Visual Sketchpad~\citep{hu2024visual} gives multimodal LMs drawing tools so reasoning proceeds over self-generated visual artifacts, improving mathematical and visual tasks.
MVoT~\citep{li2025imagine} trains MLLMs to generate image visualizations of their reasoning traces.
Our matched-modality arms serve a different purpose: they test which observation channel (text, rendered image, or both) carries the learned search controller, with corruption-based attribution (Section~\ref{sec:results}) rather than model-generated imagery.

% Paragraph purpose: Thread 4 --- runtime verification and constrained tool use.
\paragraph{Runtime verification and constrained generation.}
The LLM-Modulo framework~\citep{kambhampati2024llms} couples LLM generation with an external verifier in a generate--critique--regenerate loop, arguing that LLMs cannot self-verify.
The verifier feeds corrections back, so the loop participates in choosing the next attempt.
Our Trusted Search Runtime does not repair: it validates, applies, or rejects each operation and charges the budget, without choosing omitted operands or reordering candidates for the policy.
Self-Refine~\citep{madaan2023selfrefine} is the explicit single-model repair paradigm our runtime declines.

Constrained-decoding methods enforce structural well-formedness during generation.
PICARD~\citep{scholak2021picard} rejects inadmissible tokens at each decoding step, substantially improving text-to-SQL accuracy on Spider.
Guided generation via FSM-indexed vocabulary~\citep{willard2023efficient} provides analogous guarantees model-agnostically.
Both enforce syntactic structure without maintaining algorithmic state.
Our runtime differs in kind.
It maintains search-episode state (frontier, visited set, novelty counters, priority queue) and checks algorithm-specific invariants on every operation.
Validity and invariant compliance are reported as separate episode-level measurements, not as a single accept/reject decision on a completed output string.

% Paragraph purpose: Axis table comparing closest prior systems on the six design axes.
\begin{table}[t]
\caption{Design-axis comparison of closest prior systems.
\textbf{Algo.\ cond.}: policy conditioned on a declared algorithm's step semantics.
\textbf{Model expl.}: the model's emissions determine exploration order.
\textbf{No-repair RT}: the runtime rejects invalid operations without correction.
\textbf{Inv.\ episodes}: full episodes checked against algorithm-specific invariants.
\textbf{Match.\ mod.}: observations include non-textual modalities matched to the task.
\textbf{Transfer}: cross-domain or cross-task generalization evaluated.}
\label{tab:axis}
\centering
\footnotesize
\setlength{\tabcolsep}{3pt}
\begin{tabular}{lcccccc}
\hline
 & Algo. & Model & No-repair & Inv. & Match. & Transfer \\
\hline
Searchformer~\citep{lehnert2024beyond}   & Partial\textsuperscript{a} & Yes & No  & No  & No  & No  \\
Stream of Search~\citep{gandhi2024stream} & Partial\textsuperscript{b} & Yes & No  & No  & No  & Partial\textsuperscript{c} \\
Neural execution~\citep{velickovic2020neural} & Yes & No  & No  & No  & No  & Yes \\
Process supervision~\citep{lightman2023lets} & No  & No  & No  & Partial\textsuperscript{d} & No  & No  \\
LLM-Modulo~\citep{kambhampati2024llms}   & No  & No  & No  & Partial\textsuperscript{e} & No  & No  \\
PICARD~\citep{scholak2021picard}          & No  & Yes & Yes & No  & No  & No  \\
Sketchpad~\citep{hu2024visual}            & No  & Yes & No  & No  & Yes & No  \\
\textbf{Ours}                             & \textbf{Yes} & \textbf{Yes} & \textbf{Yes} & \textbf{Yes} & \textbf{Yes} & Null\textsuperscript{f} \\
\hline
\end{tabular}

\vspace{2pt}
{\footnotesize
\textsuperscript{a}A* dynamics serialized as tokens without typed operations or explicit operands.
\textsuperscript{b}Search streams serialized including mistakes without operation typing.
\textsuperscript{c}Self-improvement across search strategies, not cross-domain evaluation.
\textsuperscript{d}Per-step human labels for correctness, not algorithm-specific invariants.
\textsuperscript{e}External verifier provides correction feedback in a repair loop.
\textsuperscript{f}Zero-shot cross-task transfer evaluated with null result (Section~\ref{sec:results}).
}
\end{table}

% Paragraph purpose: Scoped novelty positioning derived from the axis table (D6: scoped, citation-supported, no broad priority assertion).
Table~\ref{tab:axis} compares the closest prior systems on six design axes.
Searchformer~\citep{lehnert2024beyond} is closest, training a transformer on A* search dynamics, but it operates on a single algorithm without typed operations or runtime validation.
\citet{velickovic2020neural} provide step-level algorithm imitation with cross-algorithm transfer but no runtime that checks operations during execution.
PICARD~\citep{scholak2021picard} rejects inadmissible outputs during decoding but enforces syntax on a single output string, not algorithm invariants over a stateful episode.
No prior system in the table combines all six properties.
This work combines algorithm-conditioned typed operations, a no-repair runtime with per-operation invariant checking, matched-modality observations, and cross-task transfer evaluation in a single system.
Sections~\ref{sec:design}--\ref{sec:results} report where this combination succeeds and where it does not.
